\documentclass{article}

\usepackage{float}

\begin{document}

\title{COS 314 Assignment Two Report}
% Simple multi-author: each author on its own line
\author{u24594522 -- Keegan Lewis \\
u23588579 -- Tshepiso Makelana \\
u21516261 -- Byron Pillay}
\date{April 2026}
\maketitle

\section{Introduction}

\section{Genetic Algorithm Configuration}

The Genetic Algorithm (GA) was configured with a specific set of the Knapsack Problem. The operators and details are listed below:

\begin{itemize}
    \item \textbf{Termination Criterion:} The algorithm executes for a fixed number of generations, defined by the \texttt{termination} parameter passed to the execution function.
    
    \item \textbf{Population Size:} A constant population of $N = 26$ individuals was maintained throughout the evolutionary process.
    
    \item \textbf{Initial Population and Repair:} The initial population was generated by randomly selecting approximately 25\% of the genes for each chromosome. To ensure feasibility, any individual that exceeded the knapsack capacity was subjected to a \textit{Random Repair} mechanism, where randomly selected items were removed until the total weight was within limits.
    
    \item \textbf{Selection Mechanism:} A hybrid selection strategy was employed:
    \begin{itemize}
        \item \textbf{Elitism:} The single best individual from the current population is preserved and copied directly into the next generation.
        \item \textbf{Roulette Wheel Selection:} The remaining 25 slots are filled using fitness-proportionate selection, where the probability of an individual being chosen was based on their fitness.
    \end{itemize}
    
    \item \textbf{Genetic Operators:}
    \begin{itemize}
        \item \textbf{Crossover:} A \textit{Greedy Intersection Crossover} was applied with a probability of $P_c = 0.85$. This operator preserves items common to both parents and fills the remaining capacity using a greedy approach based on the highest value-to-weight ratios.
        \item \textbf{Mutation:} Bit-flip mutation was applied with a probability of $P_m = 1/n$ (where $n$ is the number of items). 
    \end{itemize}
    
    \item \textbf{Population Update:} After the genetic operators, the population was update in a greedy style where only the best of the parents and children would move on to the next generation.
\end{itemize}

\section{Iterated Local Search Configuration}

\section{Experimental Setup}

\section{Results}

\begin{table}[H]
\centering
\begin{tabular}{|c|c|c|c|c|c|}
\hline
Problem Instance & Algorithm & Seed Value & Best Solution & Known Optimum & Runtime (seconds)\\
\hline
f1 l-d kp 10 269 & ILS &  & 295.0 &295.0 & \\
 & GA & & & & \\
\hline
f2 l-d kp 20 878 & ILS &  & 1024 & 1024.0 & \\
 & GA & & & & \\
\hline
\end{tabular}

\end{table}



\end{document}